% analysis/part_intro.tex

\chapter*{Analysis}
\phantomsection
\addcontentsline{toc}{chapter}{Analysis}

This part establishes how AIM operates as a coherent processing system
rather than as a disconnected set of models.

% ============================================================
\section*{What do we already know?}
% ============================================================

The previous parts established the conceptual, geometric, operational,
realization-aware, and modelling foundations of AIM.

Railway alignments can now be represented through sparse semantic
models, reconstructed through runtime services, interpreted through
state operators, and related to physical reality through realization
concepts.

The individual building blocks of the framework therefore exist.

However, a collection of components does not yet constitute a complete
system.

The system also has a development lineage that clarifies why these components
exist. Berlinish began as an originating research grammar for transition
functions. axtranNew was intended to turn that grammar into solvable alignment
calculations. transitionDB and TransEd address preservation and inspection of
the reusable function knowledge. SPOT and alignmentOS address concrete object
identity, state, editing, and operational use. The Knowledge Kernel provides
the stable terminology and responsibility boundaries that the originating
calculation work did not by itself supply. Research challenges and extends
claims; the Thesis explains, relates, and transfers them.

These responsibilities form an execution and knowledge chain. They must not be
read as a claim that every named component is a canonical Kernel concept or a
complete implementation.

% ============================================================
\section*{What is still missing?}
% ============================================================

Practical railway information modelling requires more than isolated
models and operators.

Data must move through processing chains.

States must be transformed.

Services must interact.

Failures must be detected.

Results must be generated from intermediate representations.

The missing element is therefore the processing architecture that
connects individual concepts into a coherent operational framework.

AIM must be understood not only as a collection of models, but as a
pipeline of interacting transformations.

% ============================================================
\section*{What comes next?}
% ============================================================

This part introduces the system-level interpretation of AIM.

The focus shifts from individual concepts toward the interaction
between concepts.

The chapters that follow investigate:
\begin{itemize}
\item transformation pipelines,
\item runtime services,
\item operator chains,
\item system interactions,
\item failure modes,
\item and architectural design principles.
\end{itemize}

The objective is to explain how railway information flows through the
AIM framework from acquisition to evaluation, realization, and
optimization.

This perspective ultimately reveals AIM as an operator-based railway
information system rather than a mere collection of geometric methods.

% ============================================================
\section*{Relation to the AIM Roadmap}
% ============================================================

Within the AIM roadmap, the present part corresponds to the transition
from alignment models toward processing pipelines and system services.

The roadmap sketch below answers the guiding question: how does the
argument progress from model structure to operational system behaviour?

\begin{center}
\begin{tikzpicture}[
  node distance=1.4cm,
  box/.style={
    draw,
    rounded corners,
    minimum width=5.0cm,
    minimum height=0.9cm,
    align=center
  }
]

\node[box] (model)
{Alignment Model};

\node[box, below=of model] (pipeline)
{Pipelines and Services};

\node[box, below=of pipeline] (system)
{System View};

\draw[->, thick] (model) -- (pipeline);
\draw[->, thick] (pipeline) -- (system);

\end{tikzpicture}
\end{center}

The progression
\[
\text{Alignment Model}
\rightarrow
\text{Pipelines and Services}
\rightarrow
\text{System View}
\]
marks the transition from static railway representations toward
operational processing architectures.

This system perspective prepares the subsequent treatment of
optimization, engineering workflows, and practical deployment.

% ============================================================
\begin{principle}[Pipeline Principle]
\label{princ:pipeline-principle}
Within AIM, railway information is processed through explicit operator
pipelines in which states, models, realizations, and evaluations are
connected by well-defined transformations.
\end{principle}

% ============================================================
\begin{summary}
Modelling explains how railway information is represented.

Analysis explains how railway information is processed.

The present part therefore establishes the system architecture that
connects models, operators, services, and transformations into a
coherent AIM framework.
\end{summary}

This system-level interpretation prepares the Optimization part, where
pipeline outputs are used for structured improvement and decision
support.
